\documentclass[11pt,paper=a4,answers]{exam}
\usepackage{graphicx,lastpage}
\usepackage{upgreek}
\usepackage{censor}
\usepackage{gensymb}
\usepackage{amsmath}
\usepackage{multicol}
\usepackage{enumitem}
\usepackage{tasks}
\usepackage{adjustbox}
\usepackage{xcolor}
\usepackage{tfrupee}
\setlength{\voffset}{0.25in}
\usepackage{tabularx}
\newcolumntype{Y}{>{\centering\arraybackslash}X}
%==============================================================
% Customise Non-Essential Packages Below
% =============================================================
\usepackage[most]{tcolorbox}
\usepackage{eso-pic}
\usepackage{tikz}
\AddToShipoutPictureBG{%
\begin{tikzpicture}[remember picture,overlay]
\foreach \x in {-8,-4,0,4,8,12,16,20}
\foreach \y in {-8,-4,0,4,8,12,16,20} {
\node[opacity=0.05,rotate=45,anchor=center] at ([xshift=\x cm,yshift=\y cm]current page.center) {%
\begin{minipage}{2cm}
\centering
\includegraphics[width=2cm]{images/logo_black.png}\\
\small preppeo.com
\end{minipage}%
};
}
\end{tikzpicture}%
}
\printanswers

%==============================================================
% Customise Non-Essential Packages Above
% =============================================================
\definecolor{svgray}{rgb}{0.90, 0.90, 0.90}
\graphicspath{{images/}}

\usepackage[normalem]{ulem}
\renewcommand\ULthickness{1pt}   %%---> For changing thickness of underline
\setlength\ULdepth{3.5ex}%\maxdimen ---> For changing depth of underline
\renewcommand{\baselinestretch}{1}
\chead{
\includegraphics[scale=0.1,height=2cm ]{logo.png}
}
\pagestyle{empty}

\pagestyle{headandfoot}
\headrule
\cfoot{W: https://preppeo.com; M: +91-8130711689}

\begin{document}
%==============================================================
% Customise Question paper details below this
% =============================================================
\begin{center}
    \centering
    \renewcommand{\arraystretch}{1.5}
    \begin{tabularx}{\textwidth}{YYY}
    & \normalsize\bf{{\Large{{Quant}}}} & \\
    By Preppeo & \normalsize\bf{{sat\_lid\_003}} & SAT\\
    \normalsize{{\textbf{{Unit:}} Algebra}} & \normalsize{{\textbf{{Chapter:}} Linear equations in one variable}} & \normalsize{{\textbf{{Topic:}} Number of Solutions (Zero, One, Infinite)}} \\
    \end{tabularx}
\end{center}
\par\noindent\rule{\textwidth}{1pt}
% =============================================================
% Customise Question paper details above this
% =============================================================

\begin{questions}
\pointsinrightmargin
\pointsdroppedatright
\marksnotpoints
\pointpoints{mark}{marks}
\pointformat{\boldmath\themarginpoints}
\bracketedpoints

% =============================================================
% Question Begin
% =============================================================
%Q: 1
% Category: Practice | Difficulty: Medium | Question Type: Multiple Choice
\question
Which of the following equations has infinitely many solutions?
\begin{tasks}(2)
    \task $4x + 5 = 4x + 10$
    \task $2(x + 3) = 2x + 6$
    \task $3x + 1 = 3x - 1$
    \task $5x + 2 = 2x + 5$
\end{tasks}

\begin{solution}
An equation has infinitely many solutions when both sides are mathematically identical. 
\\[10pt]
Let us examine Option B. By distributing the 2 on the left side, we get:
\\[10pt]
$2x + 6 = 2x + 6$ 
\\[10pt]
Because the coefficients of $x$ are equal (2 on both sides) and the constant terms are equal (6 on both sides), any value substituted for $x$ will make the equation true. This is known as an identity.

\vspace{20pt}
Answer: \colorbox{yellow!100}{B}
\end{solution}

%Q: 2
% Category: Practice | Difficulty: Medium | Question Type: Multiple Choice
\question
$ax + 12 = 5x + 12$
In the equation above, $a$ is a constant. If the equation has infinitely many solutions, what is the value of $a$?
\begin{tasks}(2)
    \task 0
    \task 5
    \task 12
    \task 17
\end{tasks}

\begin{solution}
For a linear equation to have infinitely many solutions, the coefficient of the variable and the constant term must be the same on both sides of the equal sign. 
\\[10pt]
The equation is given as $ax + 12 = 5x + 12$. 
\\[10pt]
We can see that the constant terms are already identical (both are 12). For the equation to be an identity, the coefficients of $x$ must also match. Therefore, $a$ must be equal to 5.

\vspace{20pt}
Answer: \colorbox{yellow!100}{B}
\end{solution}

%Q: 3
% Category: Practice | Difficulty: Medium | Question Type: Multiple Choice
\question
$3x + 7 = 3x + b$
In the equation above, $b$ is a constant. If the equation has no solution, which of the following could be the value of $b$?
\begin{tasks}(2)
    \task 0
    \task 7
    \task Both 0 and 7
    \task Neither 0 nor 7
\end{tasks}

\begin{solution}
An equation has no solution when the variable terms cancel out but the resulting numerical statement is false (a contradiction). 
\\[10pt]
In the equation $3x + 7 = 3x + b$, subtracting $3x$ from both sides results in $7 = b$. 
\\[10pt]
If $b = 7$, the equation becomes $7 = 7$, which means infinitely many solutions. If $b$ is any value other than 7, the statement will be false. Since 0 is not equal to 7, if $b = 0$, the equation $7 = 0$ is impossible. Thus, $b=0$ results in no solution.

\vspace{20pt}
Answer: \colorbox{yellow!100}{A}
\end{solution}

%Q: 4
% Category: Practice | Difficulty: Medium | Question Type: Multiple Choice
\question
Which of the following equations has exactly one solution?
\begin{tasks}(2)
    \task $x + x = 2x$
    \task $5(x - 1) = 5x - 5$
    \task $2x + 10 = 3x + 10$
    \task $4x - 8 = 4x + 2$
\end{tasks}

\begin{solution}
An equation has exactly one solution if the coefficients of the variable on each side of the equation are different. 
\\[10pt]
Let us analyze Option C: $2x + 10 = 3x + 10$. 
\\[10pt]
If we subtract $2x$ from both sides, we get $10 = x + 10$. Subtracting 10 from both sides yields $x = 0$. Since we found a specific value for $x$, there is exactly one solution. 
\\[10pt]
In contrast, Options A and B are identities (infinite solutions), and Option D is a contradiction (no solution).

\vspace{20pt}
Answer: \colorbox{yellow!100}{C}
\end{solution}

%Q: 5
% Category: Practice | Difficulty: Medium | Question Type: Multiple Choice
\question
$kx + 5 = 5$
In the equation above, $k$ is a constant. If the equation is true for all values of $x$, what is the value of $k$?
\begin{tasks}(2)
    \task 0
    \task 1
    \task 5
    \task Any value
\end{tasks}

\begin{solution}
The phrase "true for all values of $x$" is another way of saying the equation has infinitely many solutions. 
\\[10pt]
For the equation $kx + 5 = 5$ to be true for any $x$, the left side must simplify to exactly the same thing as the right side. 
\\[10pt]
The right side can be thought of as $0x + 5$. Therefore, the coefficient $k$ on the left side must be 0 to match the right side. If $k=0$, the equation becomes $0 + 5 = 5$, which is $5 = 5$.

\vspace{20pt}
Answer: \colorbox{yellow!100}{A}
\end{solution}

%Q: 6
% Category: Practice | Difficulty: Medium | Question Type: Multiple Choice
\question
If $3(2x + 4) = ax + b$ has infinitely many solutions, what is the value of $a + b$?
\begin{tasks}(2)
    \task 6
    \task 12
    \task 18
    \task 24
\end{tasks}

\begin{solution}
First, we distribute the 3 on the left side of the equation to see the standard form:
\\[10pt]
$6x + 12 = ax + b$ 
\\[10pt]
For the equation to have infinitely many solutions, it must be an identity where $a = 6$ and $b = 12$. 
\\[10pt]
The question asks for the sum $a + b$:
\\[10pt]
$6 + 12 = 18$

\vspace{20pt}
Answer: \colorbox{yellow!100}{C}
\end{solution}

%Q: 7
% Category: Practice | Difficulty: Medium | Question Type: Multiple Choice
\question
$4x + 1 = 4x + c$
If the equation above has no solution, which of the following is a possible value for $c$?
\begin{tasks}(2)
    \task 1
    \task 0
    \task Both 1 and 0
    \task None of the above
\end{tasks}

\begin{solution}
When an equation has no solution, the coefficients of the variable are identical but the constants are different. 
\\[10pt]
In $4x + 1 = 4x + c$, the $x$ terms will cancel out when we subtract $4x$ from both sides, leaving $1 = c$. 
\\[10pt]
If $c = 1$, we get $1 = 1$ (infinite solutions). If $c = 0$, we get $1 = 0$, which is impossible. Therefore, $c = 0$ results in no solution.

\vspace{20pt}
Answer: \colorbox{yellow!100}{B}
\end{solution}

%Q: 8
% Category: Practice | Difficulty: Medium | Question Type: Multiple Choice
\question
How many solutions does the equation $2(x - 4) = 2x - 4$ have?
\begin{tasks}(2)
    \task Zero
    \task Exactly one
    \task Infinitely many
    \task Exactly two
\end{tasks}

\begin{solution}
We begin by distributing the 2 on the left side of the equation:
\\[10pt]
$2x - 8 = 2x - 4$ 
\\[10pt]
Next, we attempt to isolate $x$ by subtracting $2x$ from both sides:
\\[10pt]
$-8 = -4$ 
\\[10pt]
Because the variable $x$ has been eliminated and we are left with a false statement ($-8$ is never equal to $-4$), there is no value of $x$ that can satisfy the equation.

\vspace{20pt}
Answer: \colorbox{yellow!100}{A}
\end{solution}

%Q: 9
% Category: Practice | Difficulty: Hard | Question Type: Multiple Choice
\question
$\frac{1}{2}(4x + 10) = ax + 5$
If the equation above has infinitely many solutions, what is the value of $a$?
\begin{tasks}(2)
    \task 2
    \task 4
    \task 5
    \task 10
\end{tasks}

\begin{solution}
We distribute the $\frac{1}{2}$ to both terms inside the parentheses:
\\[10pt]
$\frac{1}{2}(4x) + \frac{1}{2}(10) = ax + 5$ \\
$2x + 5 = ax + 5$ 
\\[10pt]
For the equation to have infinitely many solutions, it must be an identity. Comparing the two sides, the constant terms are already equal (5). The coefficients of $x$ must also be equal, so $a$ must be 2.

\vspace{20pt}
Answer: \colorbox{yellow!100}{A}
\end{solution}

%Q: 10
% Category: Practice | Difficulty: Medium | Question Type: Multiple Choice
\question
$5x - 3 = 5x - 3$
How many solutions does the given equation have?
\begin{tasks}(2)
    \task Zero
    \task Exactly one
    \task Infinitely many
    \task Exactly two
\end{tasks}

\begin{solution}
Observe that the expression on the left side is identical to the expression on the right side. 
\\[10pt]
If we add 3 to both sides, we get $5x = 5x$. If we divide by 5, we get $x = x$. This statement is true for any value of $x$ that we choose. In SAT terms, this is an identity with infinite solutions.

\vspace{20pt}
Answer: \colorbox{yellow!100}{C}
\end{solution}

%Q: 11
% Category: Practice | Difficulty: Medium | Question Type: Multiple Choice
\question
$mx + 10 = 2x + 5$
If the equation above has exactly one solution, which of the following cannot be the value of $m$?
\begin{tasks}(2)
    \task -2
    \task 0
    \task 2
    \task 10
\end{tasks}

\begin{solution}
A linear equation in the form $Ax + B = Cx + D$ has exactly one solution if and only if $A \neq C$. 
\\[10pt]
In this problem, $A = m$ and $C = 2$. If $m = 2$, the equation becomes $2x + 10 = 2x + 5$. Subtracting $2x$ would leave $10 = 5$, which has no solution. 
\\[10pt]
Therefore, for there to be exactly one solution, $m$ can be any value except 2.

\vspace{20pt}
Answer: \colorbox{yellow!100}{C}
\end{solution}

%Q: 12
% Category: Practice | Difficulty: Medium | Question Type: Multiple Choice
\question
$8x + 4 = 2(4x + 2)$
How many solutions does the given equation have?
\begin{tasks}(2)
    \task Zero
    \task Exactly one
    \task Infinitely many
    \task exactly two
\end{tasks}

\begin{solution}
Simplify the right side of the equation by distributing the 2:
\\[10pt]
$8x + 4 = 8x + 4$ 
\\[10pt]
Both sides are identical. No matter what value is used for $x$, the left side will always calculate to the same result as the right side. This characterizes an equation with infinitely many solutions.

\vspace{20pt}
Answer: \colorbox{yellow!100}{C}
\end{solution}

%Q: 13
% Category: Practice | Difficulty: Medium | Question Type: Multiple Choice
\question
Which value of $c$ makes $5x + c = 5x + 7$ have no solution?
\begin{tasks}(2)
    \task 7
    \task 0
    \task Any value except 7
    \task Any value except 0
\end{tasks}

\begin{solution}
The condition for "no solution" in a linear equation is that the variable terms cancel out while the constants remain unequal. 
\\[10pt]
In the equation $5x + c = 5x + 7$, subtracting $5x$ from both sides gives $c = 7$. 
\\[10pt]
If $c = 7$, we get $7 = 7$, which is infinite solutions. If $c$ is anything else (like 0, 5, or 10), the statement $c = 7$ will be false. Thus, any value except 7 will result in no solution.

\vspace{20pt}
Answer: \colorbox{yellow!100}{C}
\end{solution}

%Q: 14
% Category: Practice | Difficulty: Medium | Question Type: Multiple Choice
\question
$4(x + 1) = 4x + 4$
How many solutions?
\begin{tasks}(2)
    \task 0
    \task 1
    \task Infinitely many
    \task 4
\end{tasks}

\begin{solution}
Distribute the 4 on the left side:
\\[10pt]
$4x + 4 = 4x + 4$ 
\\[10pt]
Because the left side and the right side of the equation are the same, every real number is a solution.

\vspace{20pt}
Answer: \colorbox{yellow!100}{C}
\end{solution}

%Q: 15
% Category: Practice | Difficulty: Medium | Question Type: Multiple Choice
\question
$2x + 8 = 2x + 10$
How many solutions?
\begin{tasks}(2)
    \task Zero
    \task One
    \task Infinite
    \task Two
\end{tasks}

\begin{solution}
Subtract $2x$ from both sides:
\\[10pt]
$8 = 10$ 
\\[10pt]
Since 8 can never be equal to 10, there is no possible value for $x$ that makes this equation valid.

\vspace{20pt}
Answer: \colorbox{yellow!100}{A}
\end{solution}

%Q: 16
% Category: Practice | Difficulty: Medium | Question Type: Multiple Choice
\question
If $5x + a = bx + 4$ has infinitely many solutions, what is $a$?
\begin{tasks}(2)
    \task 4
    \task 5
    \task $x$
    \task 0
\end{tasks}

\begin{solution}
Infinite solutions require identical sides. 
\\[10pt]
Coefficient of $x$ on the left (5) must equal coefficient of $x$ on the right ($b$). So $b = 5$. 
\\[10pt]
The constant on the left ($a$) must equal the constant on the right (4). So $a = 4$.

\vspace{20pt}
Answer: \colorbox{yellow!100}{A}
\end{solution}

%Q: 17
% Category: Practice | Difficulty: Medium | Question Type: Multiple Choice
\question
$10x - 5 = 10x - 5$. Solutions?
\begin{tasks}(2)
    \task Zero
    \task One
    \task Infinitely many
    \task 10
\end{tasks}

\begin{solution}
The equation is an identity. Every value of $x$ satisfies the equation.

\vspace{20pt}
Answer: \colorbox{yellow!100}{C}
\end{solution}
%Q: 17
% Category: Practice | Difficulty: Hard | Question Type: Multiple Choice
\question
A storage tank has 5,000 liters. After 2 hours of draining at a constant rate, 4,200 liters remain. How many total hours to reach 2,600 liters?
\begin{tasks}(2)
    \task 4
    \task 5
    \task 6
    \task 8
\end{tasks}

\begin{solution}
To solve this problem, we must first determine the constant rate at which water is leaving the tank. 
\\[10pt]
We know that in 2 hours, the volume dropped from 5,000 to 4,200 liters. This means a total of 800 liters were removed. By dividing 800 liters by 2 hours, we find that the drainage rate is exactly 400 liters per hour. 
\\[10pt]
Next, we identify our target: reaching 2,600 liters. To find how much total water must be removed from the original 5,000 liters, we subtract: 5,000 - 2,600 = 2,400 liters total to be drained. 
\\[10pt]
Finally, we divide the total volume to be removed (2,400 liters) by the hourly rate (400 liters per hour). This calculation reveals that it will take a total of 6 hours from the start to reach the desired level.

\vspace{15pt}
Answer: \colorbox{yellow!100}{C}
\end{solution}

%Q: 18
% Category: Practice | Difficulty: Easy | Question Type: Multiple Choice
\question
What value of $p$ satisfies $7p + 140 = 700$?
\begin{tasks}(2)
    \task 80
    \task 120
    \task 560
    \task 700
\end{tasks}

\begin{solution}
Finding the value of $p$ requires us to isolate the variable by reversing the operations performed on it. 
\\[10pt]
The current expression adds 140 to the variable term. To undo this, we subtract 140 from both sides of the equation. This results in the statement that $7p$ is equal to 560. 
\\[10pt]
Now, the variable $p$ is being multiplied by 7. We use the inverse operation of division to move the 7 to the other side. By dividing 560 by 7, we determine that $p$ must be 80. 
\\[10pt]
We can verify this by checking if 7 multiplied by 80 plus 140 equals 700. Since 560 plus 140 is indeed 700, the solution is correct.

\vspace{15pt}
Answer: \colorbox{yellow!100}{A}
\end{solution}

%Q: 19
% Category: Practice | Difficulty: Medium | Question Type: Multiple Choice
\question
If $10x - 50 = 20$, what is the value of $2x - 10$?
\begin{tasks}(2)
    \task 4
    \task 10
    \task 14
    \task 20
\end{tasks}

\begin{solution}
This question tests your ability to recognize algebraic proportions rather than simply solving for a single variable. 
\\[10pt]
If we compare the left side of the equation, $10x - 50$, to the expression we are looking for, $2x - 10$, we can see that every term in the target expression is exactly one-fifth of the terms in the original equation. 
\\[10pt]
Specifically, $10x \div 5 = 2x$ and $50 \div 5 = 10$. Because of the property of equality, if we perform this division on the left side, we must also divide the right side by 5 to keep the equation balanced. 
\\[10pt]
Dividing the right side, 20, by 5 gives us 4. Therefore, without ever finding the value of $x$ itself, we can conclude that the value of $2x - 10$ is 4.

\vspace{15pt}
Answer: \colorbox{yellow!100}{A}
\end{solution}

%Q: 20
% Category: Practice | Difficulty: Easy | Question Type: Multiple Choice
\question
$15x + 5 = 155$. What is $x$?
\begin{tasks}(2)
    \task 10
    \task 11
    \task 15
    \task 150
\end{tasks}

\begin{solution}
To solve for $x$, we need to clear the constant and the coefficient away from the variable step by step. 
\\[10pt]
First, we subtract 5 from both sides of the equal sign. This leaves the variable term alone on the left: $15x = 150$. 
\\[10pt]
The next step is to isolate $x$ by removing the coefficient of 15. Since the 15 is multiplying the $x$, we divide both sides by 15. 
\\[10pt]
When we divide 150 by 15, we get a result of 10. This means that $x$ must be 10 for the equation to be balanced.

\vspace{15pt}
Answer: \colorbox{yellow!100}{A}
\end{solution}

%Q: 21
% Category: Practice | Difficulty: Medium | Question Type: Multiple Choice
\question
A painter charges 200 for supplies plus 40 per hour. If the total bill was 600, how many hours, $h$, did he work?
\begin{tasks}(2)
    \task $40h + 200 = 600$
    \task $200h + 40 = 600$
    \task $240h = 600$
    \task $40h = 600$
\end{tasks}

\begin{solution}
This problem asks us to translate a real-world billing scenario into a mathematical equation. 
\\[10pt]
The cost consists of two parts: a fixed amount for supplies that does not change, and a variable amount that depends on the duration of the job. The fixed cost is 200. The variable cost is 40 for every hour $h$, which is written as $40h$. 
\\[10pt]
By adding these two parts together, we get the total bill expression: $40h + 200$. We are told the actual bill was 600, so we set the expression equal to that value. 
\\[10pt]
To solve for $h$, we would subtract 200 from 600 to find the labor cost (400) and then divide by 40 to find the total hours (10).

\vspace{15pt}
Answer: \colorbox{yellow!100}{A}
\end{solution}

%Q: 22
% Category: Practice | Difficulty: Medium | Question Type: Multiple Choice
\question
If $4x + 2 = 10$, what is $8x - 5$?
\begin{tasks}(2)
    \task 11
    \task 15
    \task 16
    \task 20
\end{tasks}

\begin{solution}
We can find the answer by solving for $x$ and then using that value to evaluate the second expression. 
\\[10pt]
Starting with the equation $4x + 2 = 10$, we subtract 2 from both sides to get $4x = 8$. Dividing by 4 reveals that $x = 2$. 
\\[10pt]
Now that we have the value of $x$, we substitute it into the expression $8x - 5$. This becomes $8(2) - 5$. 
\\[10pt]
Following the order of operations, we multiply 8 by 2 to get 16, and then subtract 5, resulting in a final answer of 11.

\vspace{15pt}
Answer: \colorbox{yellow!100}{A}
\end{solution}

%Q: 23
% Category: Practice | Difficulty: Easy | Question Type: Multiple Choice
\question
$2x + 3x + 10 = 60$. Value of $x$?
\begin{tasks}(2)
    \task 10
    \task 12
    \task 25
    \task 50
\end{tasks}

\begin{solution}
The first priority in this equation is to combine terms that are alike. Since $2x$ and $3x$ are both terms of the same variable, they can be added together to simplify the left side to $5x + 10 = 60$. 
\\[10pt]
Next, we move the constant 10 to the other side by subtracting it from 60, which gives us $5x = 50$. 
\\[10pt]
Finally, we isolate $x$ by dividing both sides of the equation by 5. Dividing 50 by 5 shows that $x$ is 10.

\vspace{15pt}
Answer: \colorbox{yellow!100}{A}
\end{solution}

%Q: 24
% Category: Practice | Difficulty: Easy | Question Type: Multiple Choice
\question
$(w + 2) - 5 = 10$. What is $w$?
\begin{tasks}(2)
    \task 3
    \task 7
    \task 13
    \task 17
\end{tasks}

\begin{solution}
We begin by simplifying the left side of the equation. Since there is no multiplication occurring outside the parentheses, they can be removed, allowing us to combine the constants $+2$ and $-5$. 
\\[10pt]
This results in the simplified equation $w - 3 = 10$. To solve for $w$, we need to move the $-3$ to the right side of the equation. 
\\[10pt]
By using the inverse operation of addition, we add 3 to both sides. Adding 3 to 10 gives us a final value of 13 for the variable $w$.

\vspace{15pt}
Answer: \colorbox{yellow!100}{C}
\end{solution}

%Q: 25
% Category: Practice | Difficulty: Easy | Question Type: Multiple Choice
\question
$5(x - 2) = 25$. Same solution as?
\begin{tasks}(2)
    \task $x - 2 = 5$
    \task $x - 2 = 20$
    \task $5x - 2 = 25$
    \task $5x - 10 = 5$
\end{tasks}

\begin{solution}
To find an equivalent equation, we can apply an operation to both sides of the original equation that simplifies its structure. 
\\[10pt]
The original equation is $5$ multiplied by the quantity $(x - 2)$ equals 25. We can undo this multiplication by dividing both sides by 5. 
\\[10pt]
On the left side, the 5s cancel out, leaving just $x - 2$. On the right side, we divide 25 by 5, which results in 5. This gives us the new equation $x - 2 = 5$. Because we performed the same operation on both sides, the value of $x$ remains the same.

\vspace{15pt}
Answer: \colorbox{yellow!100}{A}
\end{solution}

%Q: 26
% Category: Practice | Difficulty: Medium | Question Type: Numeric Entry
\question
Solve for $x$: $0.5x + 2 = 10$.

\begin{solution}
To solve for $x$, we first isolate the term containing the variable by removing the constant 2 from the left side. 
\\[10pt]
Subtracting 2 from both sides of the equation results in $0.5x = 8$. Now, we need to eliminate the coefficient of 0.5. 
\\[10pt]
We divide both sides by 0.5. It is helpful to remember that dividing by 0.5 is mathematically the same as multiplying by 2. Therefore, multiplying 8 by 2 gives us a final result of 16.

\vspace{15pt}
Answer: \colorbox{yellow!100}{16}
\end{solution}

%Q: 27
% Category: Practice | Difficulty: Medium | Question Type: Numeric Entry
\question
$3(x + 4) = 2x - 5$. What is $x$?

\begin{solution}
We begin by distributing the 3 into the parentheses on the left side of the equation, which gives us $3x + 12 = 2x - 5$. 
\\[10pt]
Next, we want to collect all terms with the variable $x$ on one side. We subtract $2x$ from both sides, which simplifies the left side to $x + 12 = -5$. 
\\[10pt]
Finally, we isolate $x$ by subtracting 12 from both sides. Subtracting 12 from $-5$ results in $-17$, which is our final answer.

\vspace{15pt}
Answer: \colorbox{yellow!100}{-17}
\end{solution}

%Q: 28
% Category: Practice | Difficulty: Medium | Question Type: Numeric Entry
\question
$7x + 14 = 7x + k$
If the equation above has no solution, what is one possible value for $k$?

\begin{solution}
No solution occurs when the variable terms cancel out, but the constant terms remain unequal. 
\\[10pt]
In this equation, both sides have $7x$. If we subtract $7x$ from both sides, we get $14 = k$. 
\\[10pt]
For the equation to have no solution, $k$ must be any value other than 14. An easy integer choice for this would be 0. If $k = 0$, the statement $14 = 0$ is false, confirming there is no solution.

\vspace{20pt}
Answer: \colorbox{yellow!100}{0}
\end{solution}

%Q: 29
% Category: Practice | Difficulty: Hard | Question Type: Numeric Entry
\question
$\frac{1}{2}(ax + 8) = 4x + 4$
If the equation above has infinitely many solutions, what is the value of $a$?

\begin{solution}
To find the value of $a$ for infinite solutions, we must make the left side look exactly like the right side. 
\\[10pt]
First, we distribute the $\frac{1}{2}$ across the parentheses:
\\[10pt]
$\frac{1}{2}(ax) + \frac{1}{2}(8) = \frac{a}{2}x + 4$
\\[10pt]
Now we compare this to the right side: $\frac{a}{2}x + 4 = 4x + 4$. The constants match perfectly. To make the $x$ coefficients match, we set $\frac{a}{2} = 4$. Multiplying both sides by 2 gives $a = 8$.

\vspace{20pt}
Answer: \colorbox{yellow!100}{8}
\end{solution}

%Q: 30
% Category: Practice | Difficulty: Medium | Question Type: Numeric Entry
\question
$10x + 5 = 5(2x + k)$
If the equation above has infinitely many solutions, what is the value of $k$?

\begin{solution}
Infinite solutions require identical expressions on both sides. We expand the right side to compare:
\\[10pt]
$5(2x) + 5(k) = 10x + 5k$
\\[10pt]
The equation is $10x + 5 = 10x + 5k$. The variable terms already match. For the constant terms to match, we set $5 = 5k$. Dividing by 5 reveals that $k$ must be 1.

\vspace{20pt}
Answer: \colorbox{yellow!100}{1}
\end{solution}

%Q: 31
% Category: Practice | Difficulty: Hard | Question Type: Numeric Entry
\question
$4x + c = ax + 9$
If the equation above has infinitely many solutions, what is the value of $a + c$?

\begin{solution}
For infinite solutions, the coefficients of $x$ must be the same and the constants must be the same. 
\\[10pt]
From the equation, we can immediately see that $a = 4$ and $c = 9$. 
\\[10pt]
The question asks for the sum $a + c$. By adding these values together ($4 + 9$), we get a final result of 13.

\vspace{20pt}
Answer: \colorbox{yellow!100}{13}
\end{solution}

%Q: 32
% Category: Practice | Difficulty: Medium | Question Type: Numeric Entry
\question
$3x - 6 = 3x + b$
If the equation has no solution, what is the only value that $b$ cannot be?

\begin{solution}
In the equation $3x - 6 = 3x + b$, the $x$ terms will always cancel out. This leaves us with the statement $-6 = b$. 
\\[10pt]
If $b = -6$, the statement is true, meaning there are infinitely many solutions. If $b$ is any other number, the statement is false, and there is no solution. Since the question asks for the one value $b$ cannot be for there to be no solution, the answer is $-6$.

\vspace{20pt}
Answer: \colorbox{yellow!100}{-6}
\end{solution}

%Q: 33
% Category: Practice | Difficulty: Hard | Question Type: Numeric Entry
\question
$2(3x + 5) = 6x + 2k$
If the equation is an identity, what is the value of $k$?

\begin{solution}
An identity is an equation that is true for all values of $x$ (infinitely many solutions). We expand the left side:
\\[10pt]
$6x + 10 = 6x + 2k$
\\[10pt]
For this to be an identity, the constant term $2k$ must equal 10. By dividing both sides of $2k = 10$ by 2, we find that $k = 5$.

\vspace{20pt}
Answer: \colorbox{yellow!100}{5}
\end{solution}

%Q: 34
% Category: Practice | Difficulty: Medium | Question Type: Numeric Entry
\question
$ax + 4 = 12x + 4$
If the equation has exactly one solution, what is one value $a$ cannot be?

\begin{solution}
Exactly one solution occurs when the coefficients of $x$ are different. 
\\[10pt]
In this case, the coefficients are $a$ and 12. If $a = 12$, the variable terms would be identical, which would lead to infinitely many solutions (since the constants 4 and 4 also match). Therefore, for there to be exactly one solution, $a$ cannot be 12.

\vspace{20pt}
Answer: \colorbox{yellow!100}{12}
\end{solution}

%Q: 35
% Category: Practice | Difficulty: Hard | Question Type: Numeric Entry
\question
$5(x - c) = 5x - 10$
If the equation has infinitely many solutions, what is the value of $c$?

\begin{solution}
First, we distribute the 5 on the left side:
\\[10pt]
$5x - 5c = 5x - 10$
\\[10pt]
For infinite solutions, the constants must match. We set $-5c = -10$. Dividing both sides by -5 gives us $c = 2$.

\vspace{20pt}
Answer: \colorbox{yellow!100}{2}
\end{solution}

%Q: 36
% Category: Practice | Difficulty: Medium | Question Type: Numeric Entry
\question
$0x = 0$
How many solutions does this equation have? (Enter 999 for infinite)

\begin{solution}
The equation $0 = 0$ is a statement that is always true, regardless of what value is used for $x$. In algebraic terms, this means the equation is an identity and possesses infinitely many solutions.

\vspace{20pt}
Answer: \colorbox{yellow!100}{999}
\end{solution}

%Q: 37
% Category: Practice | Difficulty: Medium | Question Type: Numeric Entry
\question
$0x = 10$
How many solutions does this equation have?

\begin{solution}
The equation $0 = 10$ is a statement that is always false. No matter what number is substituted for $x$, it will be multiplied by 0 and result in 0, which can never be 10. Therefore, there are zero solutions.

\vspace{20pt}
Answer: \colorbox{yellow!100}{0}
\end{solution}

%Q: 38
% Category: Practice | Difficulty: Hard | Question Type: Numeric Entry
\question
$3(x + 2) = 3x + c$
If the equation has no solution, what is one possible value for $c$?

\begin{solution}
Expanding the left side gives $3x + 6 = 3x + c$. For no solution, the constant $c$ must be any value other than 6. 
\\[10pt]
If $c = 0$, the equation becomes $3x + 6 = 3x$, which simplifies to $6 = 0$. Since this is false, $c = 0$ is a correct possible value.

\vspace{20pt}
Answer: \colorbox{yellow!100}{0}
\end{solution}
%Q: 39
% Category: Practice | Difficulty: Medium | Question Type: Numeric Entry
\question
$x + 5 = x + 5$
How many solutions? (Enter 999 for infinite)

\begin{solution}
Since both sides of the equation are identical, the equation will be true for any real number substituted for $x$. This characterizes an identity with infinitely many solutions.

\vspace{20pt}
Answer: \colorbox{yellow!100}{999}
\end{solution}

%Q: 40
% Category: Practice | Difficulty: Hard | Question Type: Numeric Entry
\question
$ax + b = 2x + 3$
If the equation has infinitely many solutions, what is the value of $a \times b$?

\begin{solution}
For infinite solutions, the sides must be identical. This means $a$ must be 2 and $b$ must be 3. 
\\[10pt]
The question asks for the product of $a$ and $b$. Multiplying 2 and 3 results in 6.

\vspace{20pt}
Answer: \colorbox{yellow!100}{6}
\end{solution}

%Q: 41
% Category: Practice | Difficulty: Medium | Question Type: Numeric Entry
\question
$10x - 2 = 10x + 2$
How many solutions?

\begin{solution}
Subtracting $10x$ from both sides results in $-2 = 2$. This is a false statement, so there are zero solutions.

\vspace{20pt}
Answer: \colorbox{yellow!100}{0}
\end{solution}

%Q: 42
% Category: Practice | Difficulty: Medium | Question Type: Numeric Entry
\question
$x = x$
How many solutions? (Enter 999 for infinite)

\begin{solution}
This is the simplest form of an identity. It is true for every value of $x$, so it has infinitely many solutions.

\vspace{20pt}
Answer: \colorbox{yellow!100}{999}
\end{solution}

%Q: 43
% Category: Practice | Difficulty: Hard | Question Type: Numeric Entry
\question
$4(x + k) = 4x + 20$
If infinite solutions, $k = ?$

\begin{solution}
Expanding the left side gives $4x + 4k = 4x + 20$. For infinite solutions, we set $4k = 20$. Dividing by 4 gives $k = 5$.

\vspace{20pt}
Answer: \colorbox{yellow!100}{5}
\end{solution}

%Q: 44
% Category: Practice | Difficulty: Medium | Question Type: Numeric Entry
\question
$x + 1 = 2x + 1$
How many solutions?

\begin{solution}
Subtracting $x$ and 1 from both sides leaves $x = 0$. Since we found one specific value for the variable, there is exactly one solution.

\vspace{20pt}
Answer: \colorbox{yellow!100}{1}
\end{solution}

%Q: 45
% Category: Practice | Difficulty: Hard | Question Type: Numeric Entry
\question
$ax + 7 = 7x + 7$
If infinite solutions, $a = ?$

\begin{solution}
For the equation to be an identity, the coefficients of $x$ must match. Thus, $a$ must be 7.

\vspace{20pt}
Answer: \colorbox{yellow!100}{7}
\end{solution}
%Q: 46
% Category: Practice | Difficulty: Medium | Question Type: Numeric Entry
\question
$1/3(3x + 9) = x + 3$
How many solutions? (Enter 999 for infinite)

\begin{solution}
Distributing the $1/3$ gives $x + 3 = x + 3$. Since the sides are identical, there are infinitely many solutions.

\vspace{20pt}
Answer: \colorbox{yellow!100}{999}
\end{solution}

%Q: 47
% Category: Practice | Difficulty: Medium | Question Type: Numeric Entry
\question
$5x = 5x + 1$
How many solutions?

\begin{solution}
Subtracting $5x$ leaves $0 = 1$. This is false, so there are zero solutions.

\vspace{20pt}
Answer: \colorbox{yellow!100}{0}
\end{solution}
%Q: 48
% Category: Practice | Difficulty: Hard | Question Type: Multiple Choice
\question
$2(x + 5) = 2x + 10$. How many solutions?
\begin{tasks}(2)
    \task Zero
    \task Exactly one
    \task Infinitely many
    \task Exactly two
\end{tasks}

\begin{solution}
We begin by distributing the 2 on the left side of the equation: $2x + 10 = 2x + 10$. 
\\[10pt]
Notice that the left side and the right side are mathematically identical. No matter what number is used for $x$, the statement will always be true. 
\\[10pt]
In algebra, this type of equation is called an identity. Because it is always true, it has infinitely many solutions.

\vspace{15pt}
Answer: \colorbox{yellow!100}{C}
\end{solution}

%Q: 49
% Category: Practice | Difficulty: Hard | Question Type: Multiple Choice
\question
$3x + 4 = 3x + 1$. How many solutions?
\begin{tasks}(2)
    \task Zero
    \task Exactly one
    \task Infinitely many
    \task Exactly two
\end{tasks}

\begin{solution}
To analyze this equation, we attempt to isolate $x$ by subtracting $3x$ from both sides. 
\\[10pt]
This causes the variable terms to cancel out completely, leaving the numerical statement $4 = 1$. 
\\[10pt]
Since 4 can never be equal to 1, this statement is always false. This means there is no value of $x$ that can ever make the equation true, so there are zero solutions.

\vspace{15pt}
Answer: \colorbox{yellow!100}{A}
\end{solution}

%Q: 50
% Category: Practice | Difficulty: Medium | Question Type: Numeric Entry
\question
If $1/2 x + 1/4 x = 6$, what is $x$?

\begin{solution}
We start by combining the fractions on the left side. To add them, we need a common denominator, which is 4. $1/2 x$ is the same as $2/4 x$. 
\\[10pt]
Adding $2/4 x$ and $1/4 x$ results in $3/4 x$, so our equation is $3/4 x = 6$. 
\\[10pt]
To solve for $x$, we multiply both sides of the equation by 4 to clear the fraction, giving $3x = 24$. Dividing by 3 reveals that $x = 8$.

\vspace{15pt}
Answer: \colorbox{yellow!100}{8}
\end{solution}
% =============================================================
% Question End
% =============================================================

\end{questions}
\begin{center}
\rule{.5\textwidth}{1pt}
\end{center}
\end{document}
